\documentclass{beamer}
\usetheme{Madrid}
\usecolortheme{dolphin}
\definecolor{charcoal}{RGB}{54,69,79}
\setbeamercolor{structure}{fg=charcoal}
\setbeamercolor{palette primary}{bg=charcoal,fg=white}
\setbeamercolor{palette secondary}{bg=charcoal!90,fg=white}
\setbeamercolor{palette tertiary}{bg=charcoal!80,fg=white}
\setbeamercolor{palette quaternary}{bg=charcoal!70,fg=white}
\setbeamercolor{frametitle}{bg=charcoal,fg=white}
\setbeamercolor{title}{bg=charcoal,fg=white}
\usepackage{hyperref}
\usepackage{graphicx}
\usepackage{booktabs}
\usepackage{listings}
\usepackage{color}
\definecolor{dkgreen}{rgb}{0,0.6,0}
\definecolor{gray}{rgb}{0.5,0.5,0.5}
\definecolor{mauve}{rgb}{0.58,0,0.82}
\lstset{
  language=Java,
  aboveskip=3mm,
  belowskip=3mm,
  showstringspaces=false,
  columns=flexible,
  basicstyle={\small\ttfamily},
  numbers=none,
  keywordstyle=\color{blue},
  commentstyle=\color{dkgreen},
  stringstyle=\color{mauve},
  breaklines=true,
  breakatwhitespace=true,
  tabsize=2
}
\title[Exception Handling]{Exception Handling}
\subtitle{COP2250: Java Programming}
\author{Kevin Pyatt, Ph.D.}
\institute{State College of Florida \\ Pyatt Labs}
\date{Week 13}
\begin{document}
\begin{frame}
  \titlepage
\end{frame}
\begin{frame}{Today's Objectives}
\begin{itemize}
  \item Understand Java's exception hierarchy
  \item Use \texttt{try}, \texttt{catch}, and \texttt{finally} blocks
  \item Throw exceptions manually with \texttt{throw}
  \item Distinguish checked from unchecked exceptions
  \item Write multiple catch blocks correctly
\end{itemize}
\end{frame}
\begin{frame}{Exception Hierarchy}
\textbf{Every exception is an object.}
\vspace{0.3cm}
\begin{tabular}{ll}
\toprule
\textbf{Type} & \textbf{Examples} \\
\midrule
\texttt{Error} & JVM failures --- do not catch \\
\texttt{Exception} (checked) & \texttt{IOException} --- must handle \\
\texttt{RuntimeException} (unchecked) & \texttt{ArithmeticException} \\
& \texttt{NullPointerException} \\
& \texttt{ArrayIndexOutOfBoundsException} \\
& \texttt{NumberFormatException} \\
\bottomrule
\end{tabular}
\vspace{0.3cm}
\textit{Checked: compiler enforces handling. Unchecked: optional but recommended.}
\end{frame}
\begin{frame}[fragile]{try-catch-finally}
\begin{lstlisting}
try {
    int[] arr = {1, 2, 3};
    System.out.println(arr[5]); // throws
} catch (ArrayIndexOutOfBoundsException e) {
    System.out.println("Caught: " + e.getMessage());
} finally {
    System.out.println("Always runs.");
}
\end{lstlisting}
\vspace{0.3cm}
\begin{itemize}
  \item \texttt{try} --- code that might fail
  \item \texttt{catch} --- runs if exception is thrown
  \item \texttt{finally} --- runs no matter what
\end{itemize}
\vspace{0.2cm}
\textit{finally runs whether or not an exception occurred. Whether or not it was caught.}
\end{frame}
\begin{frame}[fragile]{throw}
\textbf{You can throw exceptions yourself to enforce preconditions.}
\begin{lstlisting}
public static void setAge(int age) {
    if (age < 0) {
        throw new IllegalArgumentException(
            "Age cannot be negative: " + age);
    }
    this.age = age;
}
\end{lstlisting}
\vspace{0.3cm}
\begin{tabular}{ll}
\texttt{throw} & Statement that throws an exception \\
\texttt{throws} & Declares a method may throw a checked exception \\
\end{tabular}
\vspace{0.3cm}
\textit{throw immediately stops execution of the current method.}
\end{frame}
\begin{frame}[fragile]{Multiple Catch Blocks}
\begin{lstlisting}
try {
    // code that might throw either
} catch (ArithmeticException e) {
    System.out.println("Math error: " + e.getMessage());
} catch (ArrayIndexOutOfBoundsException e) {
    System.out.println("Index error: " + e.getMessage());
} catch (Exception e) {
    System.out.println("Other: " + e.getMessage());
}
\end{lstlisting}
\vspace{0.3cm}
\begin{itemize}
  \item Java runs the first matching catch block
  \item Specific exceptions before broad ones
  \item Putting \texttt{Exception} first makes all others unreachable
\end{itemize}
\end{frame}
\begin{frame}{Key Terms}
\begin{description}
  \item[Checked exception] Must be handled; compiler enforces it
  \item[Unchecked exception] Optional to handle; \texttt{RuntimeException} subclasses
  \item[\texttt{try}] Wraps code that might throw
  \item[\texttt{catch}] Handles a specific exception type
  \item[\texttt{finally}] Always executes; used for cleanup
  \item[\texttt{throw}] Manually throws an exception
  \item[\texttt{throws}] Declares a method may propagate a checked exception
\end{description}
\end{frame}
\begin{frame}{Common Mistakes}
\begin{tabular}{ll}
\toprule
\textbf{Mistake} & \textbf{Fix} \\
\midrule
Broad catch before specific & Specific first, broad last \\
Empty catch block & At minimum print \texttt{e.getMessage()} \\
Missing finally & Add it for guaranteed cleanup \\
\texttt{throw} vs \texttt{throws} confused & throw does it; throws declares it \\
Catching \texttt{Error} & Don't \\
\bottomrule
\end{tabular}
\end{frame}
\begin{frame}{Lab --- ExceptionLab}
\textbf{Implement in order:}
\begin{enumerate}
  \item \texttt{step1()} --- catch \texttt{ArrayIndexOutOfBoundsException}
  \item \texttt{step2()} --- catch \texttt{NumberFormatException} + \texttt{finally}
  \item \texttt{step4()} --- throw \texttt{IllegalArgumentException} for negative age
  \item \texttt{step5()} --- multiple catch blocks
\end{enumerate}
\vspace{0.3cm}
\textbf{Demo Requirement:} Compile live, explain what \texttt{finally} guarantees.\\
\vspace{0.3cm}
\textbf{Next:} Assignment 12 --- Array Bounds
\end{frame}
\end{document}
